\chapter{Discussion}

The experimental results answer all five research questions, and change the
picture for two of the four hypotheses. The sections below examine why the
results took the shape they did and what they mean for practical use of the
system.

\section{Interpretation of Results}

The clearest change from the corrected results is how well adaptive holds
up. Hypothesis H1 predicted that the adaptive system would outperform every
fixed strategy on documents with diverse content. Adaptive beats both
image-only and hybrid outright on all seven metrics, and is close enough to
text-only to call four of the seven a win: list structure, table structure,
word overlap, and an effective tie on text similarity (0.894 versus 0.893).
Text-only still leads narrowly on heading structure, code blocks, and
paragraph structure. Hypothesis H1 is now largely confirmed: adaptive is not
the outright winner on every single metric, but it wins where it matters
most for structural fidelity and never meaningfully falls behind on the rest.

Hypothesis H2 predicted that text-only would fail on image-heavy or
formula-rich content, while image-only and hybrid would handle diverse
content well at a higher computational cost. The strongest evidence for the
failure half of this prediction, in an earlier version of this evaluation,
was image-only's near-total collapse on code blocks (0.250). With
postprocessing correctly credited, that score is 0.972, close to every other
strategy. The underlying bug was real --- code sometimes rendered as a
screenshot instead of a fenced block --- but it was a pipeline defect, not
evidence that vision models cannot transcribe code. What survives from H2 is
the cost half of the prediction: image-only and hybrid are still far more
expensive per page than text-only or adaptive, without a matching gain in
accuracy on this text-dominant corpus.

Hybrid performs worse than every other strategy on heading structure (0.523)
and worse than text-only and adaptive on text similarity, which hypothesis
H4 did not anticipate. Providing both the page image and the extracted text
does not help. The model receives conflicting signals (the image shows one
layout, the text extraction shows another) and appears to weight the image,
which produces the lowest heading score of any strategy. More input is not
always better. Hypothesis H4 is rejected: hybrid loses to adaptive on every
one of the seven metrics.

Hypothesis H3 is confirmed. The adaptive system processes a page in 705\,ms
on average, compared to 3{,}250\,ms for image-only and 3{,}170\,ms for
hybrid, all measured on the same rented GPU (Section~4.4). The speedup comes
directly from the classifier: most thesis pages are classified as TEXT and
handled by pymupdf4llm without any LLM call. The classifier adds a small
fixed cost per page but avoids the dominant cost, model inference, on the
majority of pages. The adaptive strategy is roughly 4.5 times cheaper than
either image-only or hybrid while matching or beating both on every metric.

\section{Comparison with Related Work}

Meuschke et al.\ \autocite{meuschke2023} benchmarked ten PDF extraction
tools across four document categories and found that no tool ranked
first on all categories. This study reaches the same conclusion at the
strategy level: no single strategy wins every page or every metric. The
adaptive design is a direct response to this finding: instead of
choosing one strategy, it classifies each page and applies the most
suitable method.

Blecher et al.\ \autocite{blecher2023} introduced Nougat, a transformer
model trained specifically on academic PDFs to produce LaTeX-formatted
Markdown. Nougat achieves high accuracy on scientific papers with dense
mathematics but requires a GPU and produces output that differs
significantly from the input layout on non-formula pages. The approach
in this thesis is complementary: it targets layout fidelity over LaTeX
reconstruction and runs on consumer hardware with a general-purpose
local model.

Subramani et al.\ \autocite{subramani2020} surveyed document layout
analysis methods and identified four recurring challenges: formula
recognition, table extraction, figure captioning, and reading-order
recovery. The results in Chapter~6 confirm that formulas and reading
order are the hardest tasks for the text extraction path, while table
extraction is where the adaptive routing adds the most value.
OmniDocBench \autocite{omnidocbench2025} recently showed that multi-modal
models outperform text-only parsers on layout-rich documents. This is
consistent with the table structure result here (adaptive 0.978 versus
text-only 0.933), but the advantage narrows on text-dominant pages.

\section{Practical Implications}
\label{sec:practical_implications}

For thesis documents specifically, the text strategy is still attractive
when speed matters most: it is 4.6 times faster than adaptive, and the two
are close enough on most metrics that the choice comes down to how much the
extra table-structure accuracy is worth for a given use case. For documents
with unknown content distribution, the adaptive strategy is the safer
choice. It beats every vision-based fixed strategy outright, comes within
noise of text-only on the metrics text-only leads, and improves
substantially on table pages.

The page classifier is the key component. Its accuracy directly
determines how much of the efficiency gain is captured. A page
incorrectly classified as IMAGE sends a text-dominant page through the
VLM, wasting inference time and introducing paraphrasing errors. The
classifier in this system uses character count, embedded image count,
vector path density, and PyMuPDF's table detector, which works well for
thesis layouts but may need adjustment for other document types.

A page with a detected table does not automatically go to the vision LLM.
The system first checks whether pymupdf4llm's own extraction already produced a valid Markdown pipe table from the text layer.
If it did, the page is treated like plain text: no image is rendered and no inference cost is paid, since pymupdf4llm reads the character stream directly and reproduces simple rows and columns correctly at no cost.
Only when that extraction fails --- typically a scanned page or a table embedded as an image --- does the page fall back to the vision LLM, which reads the visual grid structure instead of the extracted characters.
This matters for complex tables specifically: merged cells, irregular column spans, or figures embedded in cells are exactly the cases pymupdf4llm's text-layer extraction struggles with, so the vision fallback is reserved for the pages that actually need it.

The current check is presence-based rather than quality-based: it asks whether pymupdf4llm produced \emph{some} pipe table, not whether that table is \emph{correct}.
A table with merged cells can still cause pymupdf4llm to emit a technically valid but wrong pipe table --- for example, flattening a merged header into the wrong number of columns --- and the current check would accept that output without ever reaching the vision path.
A refined version of this check would inspect the extracted table's structure (column count consistency across rows, for instance) rather than just its presence, catching these silent mis-extractions before they reach the final Markdown.

Formula pages currently use the general-purpose Qwen2.5-VL-7B model.
A dedicated formula recognition tool such as pix2tex could replace
this path: it is faster, cheaper, and produces cleaner LaTeX output
for mathematical content. The adaptive
architecture supports this substitution without changes to the rest of
the pipeline: a new handler for the \texttt{FORMULA} page type is
all that is needed.

\section{Threats to Validity}
\label{sec:threats_to_validity}

Internal validity is limited by the single-annotator ground truth. The
reference files were written by the thesis author. Different readers may
format the same page differently, and the scores reflect one interpretation
of what correct Markdown looks like. Results would differ if the references
were written by a different annotator.

The metrics are a construct validity limitation. They measure structural
features (heading counts, list counts, code block counts) rather than
semantic quality. A conversion that produces the right headings at the wrong
levels would score well on text similarity but fail in practice. The heading
structure metric partially addresses this, but it counts headings rather than
verifying their level assignments.

External validity is the main open question. The corpus covers three bachelor theses --- one in German, two in English --- totalling 90 pages. Three documents is a very small sample. The rankings are close and consistent across all three (Section~\ref{sec:rq5_consistency}): adaptive leads or matches text-only on table structure in every document, and the remaining gap in text similarity is small enough on each one to read as noise rather than a real difference. Even so, a sample of three cannot establish generalisability across document types.

The results depend on what thesis documents look like. The three documents in this study follow a standard single-column layout with numbered headings, code listings, and occasional tables. That structure favours text-only extraction: the text layer is clean and PyMuPDF reproduces it accurately. Multi-column scientific papers, documents with dense formula environments, or reports with irregular table structures would break this advantage. On such documents, the vision strategies would likely close the gap or overtake text-only on more metrics.

Language and script add a second dimension. The language detection module adds a prompt hint that preserves the output language, but the classifier thresholds --- particularly the 50-character minimum for TEXT classification --- were calibrated on German and English content. Right-to-left scripts such as Arabic or Hebrew, or documents mixing CJK characters with Latin text, have not been tested. The character count threshold may behave differently on short-character scripts, misclassifying pages that should route to the vision model. Testing on a two-column scientific paper and a right-to-left script document would be the most informative next step for external validation.

The model dependency is a limitation of this evaluation, not of the system. All formal VLM results use Qwen2.5-VL-7B. The architecture is model-agnostic: the server URL and model name are parameters, so any OpenAI-compatible vision-language model can be substituted without code changes. The text-only strategy has no such dependency. Its output comes from the PDF content and the pymupdf4llm version, both stable across runs.

An initial, informal test during this thesis's revision ran the system on the same three documents with a much larger model, Qwen3.6-35B-A3B. A 35B-parameter model does not fit in the GPU memory or storage available on local hardware, so this test used a rented cloud GPU instead. That first pass, run before several pipeline bugs were fixed, surfaced a useful distinction the original hypothesis did not make: several failures that looked like model capability gaps turned out to be bugs in the pipeline instead. Figures were sometimes extracted from the wrong region of a page. Code listings were rendered as pictures instead of text. A formula sometimes appeared twice: once correctly converted to \LaTeX{} by the model, and once again as a duplicate raster image the pipeline had carried over by mistake. None of this needed a different model. It needed a fix to the extraction step that ran before the page ever reached the model.

These pipeline fixes address real output-quality problems that the seven-metric framework does not fully capture. A page where code renders as a screenshot instead of a fenced block may still register a similar code block score to one where it renders correctly, since the metric counts fence markers rather than assessing whether the page layout itself is sound.

After fixing those bugs, the model comparison was repeated formally: both models were run through the corrected pipeline and scored with the same postprocessing-aware evaluation used everywhere else in this chapter. The two models are close on most metrics. The larger model shows a genuine, specific improvement in heading structure --- 0.224 higher for hybrid and 0.141 higher for image-only --- while table structure declines slightly across all three vision-based strategies, by 0.011 to 0.034, and code block score and text similarity are effectively flat. This is a more specific finding than "a bigger model helps": it helps with the one structural element that depends most on the model inferring visual hierarchy from the page image, and does not help, or slightly hurts, on a different structural element that the classifier already routes carefully regardless of model size. This remains an informal comparison outside the main experimental design, since it reuses the same three documents rather than an independent test set, but the pattern is specific enough to be worth a controlled follow-up rather than a general claim that bigger models perform better.
